\section{结论}
尽管深度学习在智慧能源中展现出明显潜力，但其工程化应用仍面临若干关键问题。首先是数据问题。智慧能源系统中的数据往往来自不同设备、不同采样频率和不同协议，存在缺失值、异常值、噪声和时间对齐误差。如果数据预处理不到位，再复杂的模型也难以稳定发挥作用。尤其在光伏预测中，云量突变、局地遮挡和传感器故障都会显著影响模型泛化。

其次是跨场景泛化问题。很多论文在单一站点上可以取得很好的结果，但换到另一地区、另一设备或另一季节后，误差会明显增大。这说明模型可能学到的是“站点特征”，而不是更一般的能源规律。元学习、迁移学习和跨站点集成是解决这一问题的重要方向，但目前仍处于快速发展阶段。

再次是可解释性问题。能源系统属于典型的高可靠性场景，调度人员通常不会仅凭“模型误差较小”就完全信任一个黑箱系统。未来如果要将深度学习更大规模地用于调度与控制，就必须给出更明确的解释机制，例如特征贡献分析、置信区间预测、物理边界校验和异常检测联动等。

最后是预测与决策之间的脱节问题。许多研究把预测误差作为唯一目标，但真实系统更关心的是调度成本、备用容量、安全约束和用户侧体验。未来智慧能源中的高价值方向，很可能不是孤立地追求更小的 RMSE，而是把预测、优化和控制纳入统一框架中，实现“预测-决策-反馈”闭环。

深度学习正在成为智慧能源研究中的关键支撑技术，尤其在光伏和风电等新能源预测任务中表现出明显优势。通过对代表性文献的分析可以看出，LSTM、GRU、CNN-LSTM、TCN 等模型能够较好地处理强时序依赖与复杂非线性关系，在短时预测场景中具有较高应用价值。同时，领域知识融合、元学习和强化学习调度案例表明，深度学习的价值并不止于提高单一预测指标，而是能够进一步服务于储能控制、微电网运行和能源管理优化。

本文配套实验也验证了这一点：在公开 LIGHT 光伏数据集上，LSTM 在 90 分钟提前预测任务中取得了最低 RMSE，优于 Persistence 和 MLP。这说明在适当的时间粒度和预测步长下，具备序列记忆能力的深度学习模型能够更好地刻画光伏功率的时间依赖关系。

但也应看到，深度学习并不是智慧能源问题的“万能解”。如果数据质量不足、模型缺乏约束、结果难以解释，那么即使测试集误差较小，也未必能在真实系统中稳定落地。因此，本文最终的判断是：深度学习在智慧能源中的核心意义，不是把模型做得更复杂，而是利用数据驱动方法提升系统对不确定性的感知、预测和应对能力。未来真正高价值的方向，将是可解释、可迁移、可部署，并能直接支撑能源决策的深度学习方法。
